\scanpage{315}
aussi: si $X$ est normal et propre sur le \uline{trait} $S$, (avec $f_*(\uline{\mathcal O}_X)=\uline{\mathcal O}_S$ si on veut), et muni d'une section $g$ sur $S$, alors tout module inversible $\uline{L}_t$ sur $\unsure{\uline{X}_t|X_t}$ ($t$ pt générique de $S$), qui est \uline{alg. équiv. à $0$} se prolonge-t-il \inserted{en $\uline{L}$} à $X$ \inserted{et $\uline{L}_s$ est alg. équiv. à $0$} ?? C'est vrai si $\uline{\mathrm{Pic}}_{X/S}$ est représentable, ce qui \deleted{d'après Raynaud} est le cas si $X/S$ est \uline{projectif} ! Peut-être est-ce tjrs le cas \deleted{\illegible{1}} (i.e. $\uline{\mathrm{Pic}}_{X/S}$ représentable si $X\to S$ propre plat et à fibres géom. normales, voire seulement intègres...) ?
\figure{\begin{tikzpicture}[x=1cm,y=1cm]
\draw (0,1.5) .. controls (.12,1.7) and (.25,1.9) .. (.27,2.1);
\draw (1.8,1.3) .. controls (1.65,1.9) and (1.7,1.96) .. (2.25,1.96) -- (4.3,2);
\end{tikzpicture}}
\mspage{314}
